\lysh{38808}{4}

\begin{alist}
\item Το Β γράφημα περιγράφει τις πληροφορίες που αφορούν τα κύματα $P$ και $S$, γιατί είναι το μόνο στο οποίο τα κύματα ξεκινούν μαζί και το $P$ διανύει πιο γρήγορα (περίπου $19\text{sec}$ πιο γρήγορα) απόσταση περίπου $200\text{km}$.
Στο γράφημα Α δεν ξεκινούν μαζί (το $S$ έχει ξεκινήσει πιο πριν) και στο Γ γράφημα, ενώ ξεκινούν μαζί, το $S$ είναι πιο γρήγορο.

\item
\begin{rlist}
\item Κοντά σε έναν σεισμογράφο τα κύματα $P$ ταξιδεύουν με ταχύτητα $7,6 \text{km/s}$, άρα η απόσταση $y$ που διανύουν ως συνάρτηση του χρόνου $t$ εκφράζεται από τη σχέση $y = 7,6t$, που είναι η εξίσωση της ευθείας $(P)$.
Τα κύματα $S$ ταξιδεύουν με ταχύτητα $4,5 \text{km/s}$. Άρα η ευθεία $(S)$ έχει εξίσωση $y = 4,5t$.

\item Η απόσταση $d$ του σεισμογράφου από το επίκεντρο του σεισμού διανύεται από το κύμα $P$ σε χρόνο $t_P$ και από το κύμα $S$ σε χρόνο $(t_P + 19)$, αφού το $S$ θα φτάσει $19\text{sec}$ αργότερα από το $P$ στον σεισμογράφο. Άρα $d = 7,6t_P$ και $d = 4,5(t_P + 19)$.
Οπότε $7,6t_P = 4,5(t_P + 19)$, δηλαδή $3,1t_P = 85,5$ και τελικά $t_P \simeq 28\text{sec}$.
Άρα $d = 7,6 \cdot 28 \simeq 212,8\text{km}$.
Τελικά, η απόσταση $d$ του σεισμογράφου από το επίκεντρο του σεισμού είναι $d \simeq 212,8\text{km}$, ο χρόνος που χρειάστηκε το κύμα $P$ για να φτάσει στον σεισμογράφο είναι $t_P \simeq 28\text{sec}$ και ο χρόνος που χρειάστηκε το κύμα $S$ να φτάσει στον σεισμογράφο είναι $t_P + 19 = 28 + 19 = 47\text{sec}$ περίπου.
\end{rlist}
\end{alist}
